% GENERATED FILE. Edit canonical lessons, then run npm run generate:books.
% canonical-source-hash: fnv1a64:879abe5e8adc3bcc
% canonical-lessons: MR-C385-cighalne, MR-C385-uphalne, MR-C385-nivalne, MR-C385-talne-averted, MR-C385-samore-jane

\chapter{When Things Get Worse}
\label{ch:mr-when-things-get-worse}
\hlchaptermodality{385}

\begin{chapteropening}
\textbf{By the end of this chapter:} \emph{I can say that a quarrel flared up, worsened or calmed down.}

Complete the last lesson of chapter 385: I can say that a quarrel flared up, worsened or calmed down.
\end{chapteropening}

\section[\texorpdfstring{\mr{चिघळणे} (cighaḷṇe)}{चिघळणे (cighaḷṇe)}]{\mr{चिघळणे} (cighaḷṇe) — to worsen; to fester}

\label{lesson:MR-C385-cighalne}



\begin{quote}
Before the new one: say the Marathi for a tight spot, then the Marathi for a test.
\end{quote}

\subsection*{You'll want to know: \mr{चिघळणे}}
\textbf{\mr{चिघळणे}} — \emph{cighaḷṇe} — \textquotedblleft{}to worsen; to fester\textquotedblright{}.

\begin{quote}
\mr{वेळेवर} \mr{बोलणं} \mr{झालं} \mr{नाही}, \mr{म्हणून} \mr{वाद} \mr{चिघळला}.
\end{quote}

\begin{quote}
\emph{veḷevar bolṇaṁ jhālaṁ nāhī, mhaṇūn vād cighaḷlā.}
\end{quote}

\begin{quote}
Nobody talked in time, so the dispute festered.
\end{quote}

Past (he …): \mr{तो} \mr{चिघळला}, \emph{to cighaḷlā}.

Of a wound or a quarrel left alone.

\begin{cousinweb}
\textbf{\mr{चिघळणे}} (\emph{cighaḷṇe}) is a Marathi verb whose origin is not settled.
\end{cousinweb}

\begin{grammarlens}[title={What you've built}]
One more word for saying what happens if.
\end{grammarlens}

\subsection*{Your turn}
\begin{itemize}
\raggedright
  \item \emph{Say it:} \emph{cighaḷṇe}
  \item \emph{Say it:} \emph{cighaḷṇe}, once more
  \item \emph{Say it:} \emph{veḷevar bolṇaṁ jhālaṁ nāhī, mhaṇūn vād cighaḷlā.} — Nobody talked in time, so the dispute festered.
  \item \emph{Recall:} say the Marathi for a tight spot, then the Marathi for a test, then say \emph{cighaḷṇe} again
\end{itemize}

\begin{tcolorbox}[breakable,colback=teal!4,colframe=teal!35!black,title={Before you move on}]
What does \mr{चिघळणे} mean? (\textquotedblleft{}To worsen; to fester\textquotedblright{}.) Say it once more.
\end{tcolorbox}
\section[\texorpdfstring{\mr{उफाळणे} (uphāḷṇe)}{उफाळणे (uphāḷṇe)}]{\mr{उफाळणे} (uphāḷṇe) — to surge up; to flare up}

\label{lesson:MR-C385-uphalne}



\begin{quote}
Before the new one: say the Marathi for a test, then the Marathi for to worsen.
\end{quote}

\subsection*{You'll want to know: \mr{उफाळणे}}
\textbf{\mr{उफाळणे}} — \emph{uphāḷṇe} — \textquotedblleft{}to surge up; to flare up\textquotedblright{}.

\begin{quote}
\mr{जुन्या} \mr{गोष्टींवरून} \mr{पुन्हा} \mr{भांडण} \mr{उफाळलं}.
\end{quote}

\begin{quote}
\emph{junyā goṣṭīṁvarūn punhā bhāṇḍaṇ uphāḷlaṁ.}
\end{quote}

\begin{quote}
The quarrel flared up again over old matters.
\end{quote}

Past (he …): \mr{तो} \mr{उफाळला}, \emph{to uphāḷlā}.

Of flames, boiling water and feelings.

\begin{cousinweb}
\textbf{\mr{उफाळणे}} (\emph{uphāḷṇe}) is a tadbhava: CDIAL lists Marathi \emph{uphāḷṇẽ}, \textquotedblleft{}to boil or bubble up, rise (of flame)\textquotedblright{}, under Sanskrit \emph{utphalati}, \textquotedblleft{}springs open, bursts\textquotedblright{} (CDIAL 1837).
\end{cousinweb}

\begin{grammarlens}[title={What you've built}]
One more word for saying what happens if.
\end{grammarlens}

\subsection*{Your turn}
\begin{itemize}
\raggedright
  \item \emph{Say it:} \emph{uphāḷṇe}
  \item \emph{Say it:} \emph{uphāḷṇe}, once more
  \item \emph{Say it:} \emph{junyā goṣṭīṁvarūn punhā bhāṇḍaṇ uphāḷlaṁ.} — The quarrel flared up again over old matters.
  \item \emph{Recall:} say the Marathi for a test, then the Marathi for to worsen, then say \emph{uphāḷṇe} again
\end{itemize}

\begin{tcolorbox}[breakable,colback=teal!4,colframe=teal!35!black,title={Before you move on}]
What does \mr{उफाळणे} mean? (\textquotedblleft{}To surge up; to flare up\textquotedblright{}.) Say it once more.
\end{tcolorbox}
\section[\texorpdfstring{\mr{निवळणे} (nivaḷṇe)}{निवळणे (nivaḷṇe)}]{\mr{निवळणे} (nivaḷṇe) — to calm down; to clear}

\label{lesson:MR-C385-nivalne}



\begin{quote}
Before the new one: say the Marathi for to worsen, then the Marathi for to surge up.
\end{quote}

\subsection*{You'll want to know: \mr{निवळणे}}
\textbf{\mr{निवळणे}} — \emph{nivaḷṇe} — \textquotedblleft{}to calm down; to clear\textquotedblright{}.

\begin{quote}
\mr{दोन} \mr{दिवसांनी} \mr{गावातील} \mr{तणाव} \mr{निवळला}.
\end{quote}

\begin{quote}
\emph{don divsānnī gāvātīl taṇāv nivaḷlā.}
\end{quote}

\begin{quote}
After two days the tension in the village eased.
\end{quote}

Past (he …): \mr{तो} \mr{निवळला}, \emph{to nivaḷlā}.

Of muddy water as well: \mr{पाणी} \mr{निवळलं}, the water cleared.

\begin{cousinweb}
\textbf{\mr{निवळणे}} (\emph{nivaḷṇe}) is a verb built on Marathi \emph{nivaḷ}, \textquotedblleft{}pure, clear\textquotedblright{}, which CDIAL lists under Sanskrit \emph{nirmala}, \textquotedblleft{}spotless, clean\textquotedblright{} (CDIAL 7366).
\end{cousinweb}

\begin{grammarlens}[title={What you've built}]
One more word for saying what happens if.
\end{grammarlens}

\subsection*{Your turn}
\begin{itemize}
\raggedright
  \item \emph{Say it:} \emph{nivaḷṇe}
  \item \emph{Say it:} \emph{nivaḷṇe}, once more
  \item \emph{Say it:} \emph{don divsānnī gāvātīl taṇāv nivaḷlā.} — After two days the tension in the village eased.
  \item \emph{Recall:} say the Marathi for to worsen, then the Marathi for to surge up, then say \emph{nivaḷṇe} again
\end{itemize}

\begin{tcolorbox}[breakable,colback=teal!4,colframe=teal!35!black,title={Before you move on}]
What does \mr{निवळणे} mean? (\textquotedblleft{}To calm down; to clear\textquotedblright{}.) Say it once more.
\end{tcolorbox}
\section[\texorpdfstring{\mr{टळणे} (ṭaḷṇe)}{टळणे (ṭaḷṇe)}]{\mr{टळणे} (ṭaḷṇe) — to be averted}

\label{lesson:MR-C385-talne-averted}



\begin{quote}
Before the new one: say the Marathi for to surge up, then the Marathi for to calm down.
\end{quote}

\subsection*{You'll want to know: \mr{टळणे}}
\textbf{\mr{टळणे}} — \emph{ṭaḷṇe} — \textquotedblleft{}to be averted\textquotedblright{}.

\begin{quote}
\mr{थोडक्यात} \mr{अपघात} \mr{टळला}.
\end{quote}

\begin{quote}
\emph{thoḍkyāt apghāt ṭaḷlā.}
\end{quote}

\begin{quote}
An accident was narrowly averted.
\end{quote}

Past (he …): \mr{तो} \mr{टळला}, \emph{to ṭaḷlā}.

Intransitive: the danger is the subject; \mr{टाळणे}, to avoid, is its causative partner.

\begin{cousinweb}
\textbf{\mr{टळणे}} (\emph{ṭaḷṇe}) is inherited: CDIAL lists Marathi \emph{ṭaḷṇẽ}, \textquotedblleft{}to pass off\textquotedblright{}, under \emph{ṭalati}, \textquotedblleft{}moves aside\textquotedblright{} (CDIAL 5450), beside \emph{ṭāḷṇẽ}, \textquotedblleft{}to avoid\textquotedblright{}.
\end{cousinweb}

\begin{grammarlens}[title={What you've built}]
One more word for saying what happens if.
\end{grammarlens}

\subsection*{Your turn}
\begin{itemize}
\raggedright
  \item \emph{Say it:} \emph{ṭaḷṇe}
  \item \emph{Say it:} \emph{ṭaḷṇe}, once more
  \item \emph{Say it:} \emph{thoḍkyāt apghāt ṭaḷlā.} — An accident was narrowly averted.
  \item \emph{Recall:} say the Marathi for to surge up, then the Marathi for to calm down, then say \emph{ṭaḷṇe} again
\end{itemize}

\begin{tcolorbox}[breakable,colback=teal!4,colframe=teal!35!black,title={Before you move on}]
What does \mr{टळणे} mean? (\textquotedblleft{}To be averted\textquotedblright{}.) Say it once more.
\end{tcolorbox}
\section[\texorpdfstring{\mr{सामोरे} \mr{जाणे} (sāmore jāṇe)}{सामोरे जाणे (sāmore jāṇe)}]{\mr{सामोरे} \mr{जाणे} (sāmore jāṇe) — to confront; to face}

\label{lesson:MR-C385-samore-jane}



\begin{quote}
Before the new one: say the Marathi for to calm down, then the Marathi for to be averted.
\end{quote}

\subsection*{You'll want to know: \mr{सामोरे} \mr{जाणे}}
\textbf{\mr{सामोरे} \mr{जाणे}} — \emph{sāmore jāṇe} — \textquotedblleft{}to confront; to face\textquotedblright{}.

\begin{quote}
\mr{ती} \mr{अडचणींना} \mr{धैर्याने} \mr{सामोरी} \mr{गेली}.
\end{quote}

\begin{quote}
\emph{tī aḍcaṇīnnā dhairyāne sāmorī gelī.}
\end{quote}

\begin{quote}
She faced the difficulties bravely.
\end{quote}

Past (he …): \mr{तो} \mr{सामोरा} \mr{गेला}, \emph{to sāmorā gelā}.

What is faced takes -\mr{ला}, and \mr{सामोरा} agrees with the subject: \mr{तो} \mr{सामोरा} \mr{गेला}, \mr{ती} \mr{सामोरी} \mr{गेली}.

\begin{cousinweb}
\textbf{\mr{सामोरे} \mr{जाणे}} (\emph{sāmore jāṇe}) joins \mr{सामोरे}, \textquotedblleft{}in front\textquotedblright{}, inherited from Sanskrit \emph{saṁmukha}, \textquotedblleft{}facing\textquotedblright{} (CDIAL 12982 lists Marathi \emph{sāmorā}), to \mr{जाणे}, \textquotedblleft{}to go\textquotedblright{}: to go up face to face.
\end{cousinweb}

\begin{grammarlens}[title={What you've built}]
One more word for saying what happens if.
\end{grammarlens}

\subsection*{Your turn}
\begin{itemize}
\raggedright
  \item \emph{Say it:} \emph{sāmore jāṇe}
  \item \emph{Say it:} \emph{sāmore jāṇe}, once more
  \item \emph{Say it:} \emph{tī aḍcaṇīnnā dhairyāne sāmorī gelī.} — She faced the difficulties bravely.
  \item \emph{Recall:} say the Marathi for to calm down, then the Marathi for to be averted, then say \emph{sāmore jāṇe} again
\end{itemize}

\begin{tcolorbox}[breakable,colback=teal!4,colframe=teal!35!black,title={Before you move on}]
What does \mr{सामोरे} \mr{जाणे} mean? (\textquotedblleft{}To confront; to face\textquotedblright{}.) Say it once more.
\end{tcolorbox}
